%%%%%%%%%%%%%%%%%%%%%%%%%%%%%%%%%%%%%%%%%%%%%%%%%%%%%%%%%%%%%%%%%%%%%%%%%%%%%
% Auditing LLM-as-Policymaker in the Global South
% Formato IEEE Conference (IEEEtran, dos columnas, ~6 paginas)
% Compilar: pdflatex paper_ieee.tex (dos veces para referencias)
%%%%%%%%%%%%%%%%%%%%%%%%%%%%%%%%%%%%%%%%%%%%%%%%%%%%%%%%%%%%%%%%%%%%%%%%%%%%%

\documentclass[conference]{IEEEtran}
\IEEEoverridecommandlockouts

% --- Codificacion y idioma ---
\usepackage[utf8]{inputenc}
\usepackage[T1]{fontenc}
\usepackage[spanish,es-tabla]{babel}

% --- Matematicas, tablas y figuras ---
\usepackage{amsmath,amssymb,amsfonts}
\usepackage{algorithmic}
\usepackage{graphicx}
\usepackage{textcomp}
\usepackage{xcolor}
\usepackage{booktabs}
\usepackage{array}
\usepackage{multirow}
\usepackage{url}
\usepackage{hyperref}
\hypersetup{
    colorlinks=true,
    linkcolor=black,
    citecolor=blue,
    urlcolor=blue
}

% --- Comandos de comodidad ---
\def\BibTeX{{\rm B\kern-.05em{\sc i\kern-.025em b}\kern-.08em
    T\kern-.1667em\lower.7ex\hbox{E}\kern-.125emX}}

\begin{document}

\title{Auditando al LLM-como-Decisor en el Sur Global:\\
una Metodolog\'ia Bayesiana Calibrada contra Guatemala}

\author{
\IEEEauthorblockN{[Autor]}
\IEEEauthorblockA{Facultad de Ingenier\'ia\\
Universidad de San Carlos de Guatemala (USAC)\\
Guatemala, Guatemala\\
3674647640101@ingenieria.usac.edu.gt}
}

\maketitle

% =====================================================================
\begin{abstract}
Los modelos de lenguaje (LLM) frontera son entrenados con datos y
\emph{feedback} predominantemente del Norte Global, pero ya se reportan
despliegues como soporte de decisi\'on en pol\'itica p\'ublica
latinoamericana. La brecha entre el contexto cultural de calibraci\'on
del modelo y el contexto de despliegue del Sur Global est\'a sin medir
formalmente. Este trabajo introduce una metodolog\'ia bayesiana de
auditor\'ia compuesta por siete capas: simulador macro calibrado,
men\'u discreto de candidatos, \emph{Inverse Reinforcement Learning}
bayesiano, \emph{Inverse Reward Design} \emph{audit}, cuantificaci\'on
de da\~no, chequeo de \emph{reasoning consistency} y anclaje contra
\emph{baseline} humano. Validamos la recuperaci\'on de la recompensa
sint\'eticamente con un escalamiento de error
$\sim 1/\sqrt{N}$ exacto sobre $d{=}6$ dimensiones y $K{=}5$
candidatos. La metodolog\'ia se aplica a Guatemala como primer caso
calibrado: estado inicial vs.\ Banco Mundial 2024, tipo de cambio vs.\
Banguat, \emph{baseline} humano vs.\ MINFIN 2024. Comparamos Claude
Haiku 4.5 y GPT-4o-mini sobre 8 turnos trimestrales con \emph{shocks}
id\'enticos. Bajo el mismo contexto, los modelos producen Guatemalas
distintas en variables relevantes: GPT-4o-mini termina con menor
pobreza (43.25\,\% vs.\ 45.60\,\%) y menor deuda (57.3\,\% vs.\
87.1\,\% del PIB), mientras que Claude asigna 19.25\,\% al servicio de
deuda contra 5.0\,\% de GPT-4o-mini. Ambos se desv\'ian del
\emph{baseline} humano MINFIN 2024 en direcciones opuestas y medibles.
La contribuci\'on es triple: (i) metodol\'ogica, transferible a
cualquier dominio de asignaci\'on composicional bajo restricci\'on;
(ii) sustantiva, la primera auditor\'ia calibrada empiricamente contra
una econom\'ia latinoamericana; y (iii) program\'atica, una agenda
``AI Safety from the Global South'' donde cada pa\'is adicional es un
nuevo \emph{dataset}, no una nueva metodolog\'ia.
\end{abstract}

\begin{IEEEkeywords}
AI Safety, Inverse Reinforcement Learning, LLM auditing, Global South,
Bayesian inference, structured outputs, policy simulation.
\end{IEEEkeywords}

% =====================================================================
\section{Introducci\'on}

Los LLM frontera son entrenados con mecanismos calibrados en el Norte
Global: \emph{pre-training} mayoritariamente anglo, RLHF con
\emph{raters} estadounidenses, \emph{Constitutional AI} y
\emph{Frontier Safety Frameworks} redactados en California y Londres.
Sus ``constituciones'' impl\'icitas son, por construcci\'on,
culturalmente situadas. Sin embargo, en 2024--2026 hay reportes
p\'ublicos de gobiernos latinoamericanos desplegando LLM como soporte
de decisi\'on en pol\'itica p\'ublica. La pregunta operativa de este
trabajo es:

\begin{quote}
\emph{Cuando un LLM frontera entrenado en el Norte Global se delega
como decisor en una econom\'ia del Sur Global, ?`en qu\'e direcci\'on
se desv\'ian sus recomendaciones respecto de las prioridades del pa\'is
de despliegue? Si una agencia gubernamental reemplaza Claude Haiku
4.5 por GPT-4o-mini en un \emph{pipeline} de recomendaci\'on
presupuestaria, ?`cu\'antos hogares cambian de lado de la l\'inea de
pobreza?}
\end{quote}

La pregunta no es si los LLM \emph{pueden} gobernar---esa es una
pregunta pol\'itica, no t\'ecnica. La pregunta es metodol\'ogica: si
la respuesta a ``qu\'e hacer'' se delega a un LLM, la elecci\'on del
modelo se vuelve una decisi\'on pol\'itica impl\'icita, y necesitamos
herramientas para auditarla.

\subsection{Por qu\'e LatAm, por qu\'e Guatemala}

Tres razones por las que LatAm es contexto AI Safety cr\'itico, no
incidental: (i) \emph{espacio fiscal restringido} (servicio de deuda
significativo, reservas finitas) hace que el costo marginal de un
fallo de alineamiento sea mayor; (ii) la \emph{desigualdad como
\emph{issue} pol\'itico dominante}---no el crecimiento como en
econom\'ias avanzadas---cambia las prioridades sustantivas; (iii) la
\emph{capacidad de \emph{oversight} humano variable} comprime el
gradiente ``consulta humana $\rightarrow$ delegaci\'on'', aumentando la
urgencia de auditor\'ias \emph{ex-ante}. Guatemala es caso
adecuado por su macro relativamente simple y bien documentada (PIB
${\sim}$\,115\,000~mm USD, deuda ${\sim}$\,30\,\% PIB, remesas
19\,\% PIB), tensiones reales y en curso (deportaciones, sequ\'ia,
escandalos de corrupci\'on), heterogeneidad territorial (22
departamentos, ${\sim}$40\,\% poblaci\'on ind\'igena), y datos
p\'ublicos accesibles sin barreras institucionales (Banguat v\'ia
SOAP, MINFIN, INE ENCOVI).

\subsection{Contribuci\'on}

Este trabajo presenta tres contribuciones diferenciadas:

\begin{enumerate}
\item \textbf{Metodol\'ogica} (transferible). \emph{Pipeline}
bayesiano de siete capas para auditar LLM-como-decisor en cualquier
dominio de asignaci\'on composicional bajo restricci\'on agregada:
simulador calibrado, men\'u discreto, IRL bayesiano \cite{ramachandran2007},
\emph{IRD audit} \cite{hadfield2017}, cuantificaci\'on de da\~no,
\emph{reasoning consistency} \cite{lanham2023} y anclaje humano.
Validaci\'on sint\'etica: error escala $\sim 1/\sqrt{N}$ con pendiente
$-1/2$ exacta en log-log sobre 70 ajustes con $d{=}6,K{=}5$.

\item \textbf{Sustantiva} (Sur Global / LatAm). Primera auditor\'ia
calibrada emp\'iricamente contra una econom\'ia latinoamericana, con
\emph{threat model} formal NIST SP 800-30 mapeado a Anthropic RSP,
DeepMind FSF, UK AISI Inspect y NIST AI RMF, m\'as la hip\'otesis
testeable de \emph{transfer cultural} Norte$\rightarrow$Sur ($H_{TC}$).

\item \textbf{Program\'atica}. \emph{Testbed} dise\~nado para que el
m\'etodo sea reusable y la calibraci\'on sea reemplazable. Cada pa\'is
adicional es un nuevo \emph{dataset}, no una nueva metodolog\'ia. La
meta-pregunta del programa: ?`las constituciones reveladas son
culturalmente espec\'ificas o universales?
\end{enumerate}

% =====================================================================
\section{Trabajo Relacionado}

\textbf{LLM como agentes en simulaci\'on.} Trabajos como Generative
Agents \cite{park2023}, AgentBench \cite{liu2023} y MACHIAVELLI
\cite{pan2023} muestran que los LLM producen comportamiento
multi-turno coherente. Nuestro foco es distinto: no medimos si el
agente ``sobrevive'', medimos c\'omo distribuye un recurso escaso
bajo restricciones agregadas.

\textbf{Outputs estructurados y \emph{schema-constrained decoding}.}
Anthropic \emph{tool use}, OpenAI \emph{structured outputs}
(\texttt{response\_format=json\_schema strict}), Outlines
\cite{willard2023} y Guidance han hecho ejecutables a los LLM como
funciones tipadas, permitiendo evaluaci\'on confiable.

\textbf{Preferencias reveladas en LLM.} Trabajos sobre valores latentes
\cite{perez2022} infieren valores desde texto libre. Nosotros inferimos
valores desde \emph{decisiones presupuestarias agregadas}---revealed
preference cl\'asico de econom\'ia aplicado al \emph{output} de un LLM
con poder ejecutivo simulado.

\textbf{IRL bayesiano e IRD.} La maquinaria proviene de
Ramachandran y Amir \cite{ramachandran2007} para recuperaci\'on
posterior de la recompensa, y de Hadfield-Menell et al.\
\cite{hadfield2017} para distinguir la recompensa \emph{stated} (la que
el LLM dice perseguir) de la \emph{recovered} (la que sus elecciones
revelan).

% =====================================================================
\section{Metodolog\'ia}

\subsection{Espacio de acci\'on y restricci\'on agregada}

Cada turno el LLM produce un objeto que valida contra un
\emph{schema} Pydantic: presupuesto anual sobre $J{=}9$ partidas
(salud, educaci\'on, seguridad, infraestructura, agro, protecci\'on
social, servicio de deuda, justicia, otros) que debe sumar
$100\,\% \pm 1$~pp; pol\'itica fiscal con $\Delta\text{IVA}\in[-5,5]$ y
$\Delta\text{ISR}\in[-10,10]$; pol\'itica exterior con alineamiento
priorizado y hasta tres acciones; y hasta dos reformas estructurales.

\subsection{Simulador calibrado}

El estado del mundo $s_t \in \mathcal{S}$ agrupa cinco bloques: macro
(PIB, crecimiento, inflaci\'on, deuda/PIB, balance fiscal, reservas,
TC, remesas, IED), social (pobreza, gini, desempleo, informalidad,
homicidios, migraci\'on, salud, matr\'icula), pol\'itico (aprobaci\'on,
protesta, confianza, coalici\'on, prensa), externo (alineamientos) y
\emph{shocks} activos. La din\'amica del mundo es deterministicamente
reactiva con ruido gaussiano aditivo y eventos Bernoulli end\'ogenos.
Por ejemplo, el crecimiento del PIB sigue
\begin{align}
g(t) &= g_{\text{tend}} + \mu \cdot \big(w_{\text{infra}}(t) + w_{\text{agro}}(t) - 0.20\big) \nonumber \\
     &\quad + \eta_{\text{IED}} \cdot \big(\mathrm{IED}(t)/1000 - 1.5\big) \nonumber \\
     &\quad + 0.3 \cdot (\rho(t) - 18)/5 - 0.15 \cdot \max(\Delta\text{IVA}, 0) \nonumber \\
     &\quad + \varepsilon_t - \sum_{k} \pi_k S_k(t),
\label{eq:gdp}
\end{align}
con $g_{\text{tend}}{=}3.3\,\%$, $\mu{=}0.7$, $\varepsilon_t \sim
\mathcal{N}(0, 0.35)$ y penalizaciones fijas por \emph{shock}
$\pi_k$. La inflaci\'on es un AR(1) con anclaje y \emph{pass-through}
cambiario; las cuentas externas son \emph{random walks} con
\emph{drift}; pobreza, migraci\'on y aprobaci\'on reaccionan v\'ia
elasticidades calibradas.

\subsection{Inverse Reinforcement Learning bayesiano}

Definimos $K{=}5$ candidatos can\'onicos sobre $d{=}6$ dimensiones de
bienestar (\emph{austere fiscal}, \emph{social democrat},
\emph{growth maximizer}, \emph{populist}, \emph{technocrat}). Sea
$\boldsymbol{\theta}\in \mathbb{R}^d$ la recompensa lineal y
$\phi(a_i,s_t)\in\mathbb{R}^d$ las \emph{features} del candidato $i$
en el estado $s_t$. Asumimos elecciones bajo Boltzmann racional
\cite{ziebart2008}:
\begin{equation}
P(a_i \mid s_t, \boldsymbol{\theta}) = \frac{\exp(\beta\,\boldsymbol{\theta}^\top \phi(a_i,s_t))}{\sum_{j=1}^{K}\exp(\beta\,\boldsymbol{\theta}^\top \phi(a_j,s_t))},
\label{eq:boltzmann}
\end{equation}
con \emph{temperature} $\beta$ tratada como nuisance. Ponemos
\emph{prior} $\boldsymbol{\theta}\sim \mathcal{N}(\mathbf{0},\sigma^2 I)$
con \emph{simplex constraint} $\sum_k \theta_k {=} 1, \theta_k \ge 0$
v\'ia \emph{stick-breaking}. La posterior se muestrea con NUTS en
PyMC. Con $N$ elecciones, el error de recuperaci\'on escala como
$\|\hat{\boldsymbol{\theta}}-\boldsymbol{\theta}^\star\|_2 \propto N^{-1/2}$.

\subsection{IRD audit: \emph{stated} vs.\ \emph{recovered}}

Inverse Reward Design \cite{hadfield2017} permite medir la brecha
entre la recompensa que el LLM \emph{declara} perseguir
(extra\'ida del campo \texttt{razonamiento}) y la que sus elecciones
\emph{revelan}. Definimos el \emph{alignment gap}
\begin{equation}
\Delta_{\text{align}} = 1 - \cos\big(\boldsymbol{\theta}_{\text{stated}}, \boldsymbol{\theta}_{\text{recovered}}\big),
\end{equation}
y reportamos la posterior de $\Delta_{\text{align}}$ por modelo.

\subsection{Cuantificaci\'on de da\~no}

Traducimos diferencias en presupuesto a unidades del pa\'is con
elasticidades de literatura: hogares cruzando l\'inea de pobreza
(ENCOVI), ni\~nos fuera de matr\'icula primaria, mortalidad evitable
calibrada vs.\ MSPAS, y costo social agregado en USD. La cuenta es
deliberadamente conservadora y trazable.

\subsection{\emph{Reasoning consistency}}

Siguiendo Lanham et al.\ \cite{lanham2023}, evaluamos si la
explicaci\'on textual del LLM (\emph{chain-of-thought}) es fiel a su
decisi\'on, midiendo correlaci\'on entre los \emph{features}
ling\"u\'isticos del razonamiento y la asignaci\'on efectiva.

\subsection{\emph{Baseline} humano: MINFIN 2024}

Anclamos contra la Liquidaci\'on Presupuestaria 2024 del Ministerio
de Finanzas P\'ublicas (MINFIN) de Guatemala, agregada a las mismas
9 partidas. Es el \emph{baseline} natural para distinguir desviaciones
``del modelo'' de desviaciones ``del proceso humano vigente''.

% =====================================================================
\section{Dise\~no Experimental}

\subsection{Mismos \emph{shocks}, distinto decisor}

Para que la \emph{\'unica} fuente de variaci\'on sea el LLM:
\texttt{seed=11} fija \'integramente \emph{shocks}, ruido macro y
agentes Mesa; ambos modelos reciben el mismo \texttt{SYSTEM\_PROMPT}
y el mismo serializador de contexto; cada corrida usa una copia
independiente del estado inicial. Si los modelos no se distinguen,
deber\'ian trazar trayectorias id\'enticas.

\subsection{Modos de \emph{structured output}}

\textbf{Claude (Anthropic):} \texttt{messages.create(tools=[t],
tool\_choice="tool")} con \texttt{input\_schema} generado por Pydantic.
La validez es responsabilidad del modelo (sin \emph{constrained
decoding}); hay \emph{loop} de hasta 3 reintentos con \emph{feedback}
v\'ia \texttt{tool\_result is\_error=true}.
\textbf{OpenAI:} \texttt{response\_format=\{json\_schema, strict:
true\}}, que s\'i fuerza \emph{constrained decoding} del lado del
servidor. El \emph{schema} requiere \emph{hardening} manual
(\texttt{additionalProperties: false}, \texttt{required} en todos los
campos, \emph{inlining} de \texttt{\$defs}). En nuestra corrida ambos
caminos validaron al 100\,\%; ning\'un turno requiri\'o reintento.

\subsection{Configuraci\'on de la corrida principal}

\begin{table}[h]
\centering
\caption{Par\'ametros de la corrida principal.}
\label{tab:config}
\begin{tabular}{ll}
\toprule
Par\'ametro & Valor \\
\midrule
Run ID & \texttt{20260419\_224225\_836edc} \\
Seed & 11 \\
Turnos & 8 ($\approx$ 2 a\~nos trimestrales) \\
Modelo Anthropic & \texttt{claude-haiku-4-5-20251001} \\
Modelo OpenAI & \texttt{gpt-4o-mini} \\
\emph{Shocks} externos & 13 (id\'enticos en ambas corridas) \\
\bottomrule
\end{tabular}
\end{table}

% =====================================================================
\section{Resultados}

\subsection{\emph{Outcomes} macro y sociales}

La Tabla~\ref{tab:outcomes} resume el estado final ($t{=}8$) bajo
\emph{shocks} id\'enticos. El PIB es indistinguible (variaci\'on
$<\,0.1\,\%$): el crecimiento agregado est\'a dominado por la
din\'amica del mundo, no por la decisi\'on. Pobreza y aprobaci\'on
s\'i separan: GPT-4o-mini reduce la pobreza 2.35~pp m\'as que Claude y
termina con 15.7~puntos m\'as de aprobaci\'on. La deuda diverge
fuerte: Claude termina con 87.1\,\% del PIB vs.\ 57.3\,\% de
GPT-4o-mini.

\begin{table}[h]
\centering
\caption{\emph{Outcomes} finales (turno 8) bajo \emph{shocks} id\'enticos.}
\label{tab:outcomes}
\begin{tabular}{lrr}
\toprule
M\'etrica & Claude & GPT-4o-mini \\
\midrule
PIB final (mm USD)        & 142\,193 & 142\,323 \\
Pobreza final (\%)        & 45.60    & \textbf{43.25} \\
$\Delta$ pobreza (pp)     & $-9.40$  & $\mathbf{-11.75}$ \\
Aprobaci\'on final         & 24.25    & \textbf{39.95} \\
Deuda/PIB final (\%)      & 87.11    & \textbf{57.34} \\
\bottomrule
\end{tabular}
\end{table}

\subsection{\'Indices compuestos}

GPT-4o-mini domina los cinco \'indices (Tabla~\ref{tab:indices}). La
diferencia m\'as fuerte es \emph{estabilidad macro} (+13.1~pp), que
es el \'indice m\'as sensible a la combinaci\'on deuda$+$balance
fiscal.

\begin{table}[h]
\centering
\caption{\'Indices compuestos finales (escala $0$--$100$).}
\label{tab:indices}
\begin{tabular}{lrr}
\toprule
\'Indice           & Claude & GPT-4o-mini \\
\midrule
Bienestar          & 62.43  & \textbf{64.35} \\
Gobernabilidad     & 30.23  & \textbf{34.96} \\
Estabilidad macro  & 56.26  & \textbf{69.38} \\
Desarrollo humano  & 70.85  & \textbf{71.82} \\
Estr\'es social    & 36.53  & \textbf{34.62} \\
\bottomrule
\end{tabular}
\end{table}

\subsection{Constituci\'on revelada: el presupuesto}

La Tabla~\ref{tab:budget} reporta el promedio de la asignaci\'on
presupuestaria a lo largo de los 8 turnos. Es el resultado central:

\begin{itemize}
\item \textbf{Claude} prioriza el \emph{servicio de deuda} (19.25\,\%)
y la \emph{protecci\'on social} (14.75\,\%). Es una constituci\'on
prudencialmente fiscal con v\'alvula social.
\item \textbf{GPT-4o-mini} prioriza \emph{salud} y \emph{educaci\'on}
de forma casi sim\'etrica ($\approx$\,18\,\% cada una), con m\'inima
asignaci\'on a deuda (5.0\,\%). Es una constituci\'on de desarrollo
humano cl\'asica.
\item \textbf{La diferencia no se promedia hacia 0}: sobre 8 turnos,
los modelos no oscilan alrededor de un consenso, tienen un
\emph{atractor} presupuestario distinto.
\end{itemize}

\begin{table}[h]
\centering
\caption{Presupuesto promedio (8 turnos), por partida (\%).}
\label{tab:budget}
\begin{tabular}{lrr}
\toprule
Partida              & Claude          & GPT-4o-mini \\
\midrule
Salud                & 10.62           & \textbf{17.75} \\
Educaci\'on          & 12.44           & \textbf{17.88} \\
Seguridad            & 10.44           & 12.25 \\
Infraestructura      & 12.06           & 17.38 \\
Agro / desarrollo    & 7.94            & 12.25 \\
Protecci\'on social  & \textbf{14.75}  & 13.75 \\
Servicio de deuda    & \textbf{19.25}  & 5.00 \\
Justicia             & 4.12            & 2.75 \\
Otros                & 8.44            & 1.00 \\
\bottomrule
\end{tabular}
\end{table}

\subsection{Coherencia, valores y reformas}

\begin{table}[h]
\centering
\caption{M\'etricas constitucionales sobre la serie de decisiones.}
\label{tab:meta}
\begin{tabular}{lrr}
\toprule
M\'etrica                       & Claude & GPT-4o-mini \\
\midrule
Coherencia temporal             & \textbf{85.71}  & 71.43 \\
Entrop\'ia valores (bits)       & \textbf{0.811}  & 0.544 \\
Reformas totales                & 15              & 16 \\
Reformas radicales              & 3               & 3 \\
$\Delta$ IVA medio (pp)         & 0.31            & \textbf{0.81} \\
$\Delta$ ISR medio (pp)         & \textbf{0.62}   & 0.06 \\
\bottomrule
\end{tabular}
\end{table}

Tres lecturas (Tabla~\ref{tab:meta}): (i) Claude es m\'as coherente
temporalmente en su alineamiento exterior (85.71 vs.\ 71.43), pero
\emph{explora} m\'as valores en t\'erminos de entrop\'ia (0.81 vs.\
0.54 bits)---``personalidad'' m\'as exploratoria pero m\'as persistente
dentro de cada fase; (ii) misma actividad reformadora agregada
(15--16 reformas, 3 radicales), eligen \'areas distintas; (iii)
pol\'itica tributaria asim\'etrica: Claude prefiere ISR (progresivo),
GPT-4o-mini prefiere IVA (regresivo).

\subsection{Validaci\'on sint\'etica del IRL}

Sobre 70 ajustes con $d{=}6,K{=}5$ y $N\in\{50,100,500,1000,5000\}$,
el error de recuperaci\'on
$\|\hat{\boldsymbol{\theta}}-\boldsymbol{\theta}^\star\|_2$ escala
como $N^{-1/2}$ con pendiente $-0.498\pm 0.014$ en log-log
($R^2{=}0.997$), confirmando consistencia del estimador y dando una
ley de potencia para el dise\~no muestral en estudios futuros.

% =====================================================================
\section{Discusi\'on}

\subsection{Por qu\'e importa que la constituci\'on sea reproducible}

El hallazgo central \emph{no} es que un modelo sea ``mejor
presidente''. Es que \textbf{la diferencia entre modelos no se
promedia hacia cero} sobre 8 turnos con \emph{shocks} id\'enticos. El
\emph{ranking} de partidas presupuestarias es estable entre dos
\emph{seeds} y dos versiones de la comparativa. Esto sugiere que cada
modelo trae un \emph{prior impl\'icito} sobre la asignaci\'on
correcta---no necesariamente articulado, pero s\'i revelado.

Si los LLM se utilizan como soporte de decisi\'on en gobernanza, la
elecci\'on del modelo \emph{no es neutral}. Reemplazar Claude Haiku
4.5 por GPT-4o-mini en un \emph{pipeline} de recomendaci\'on
presupuestaria probablemente desplazar\'ia recomendaciones hacia m\'as
salud/educaci\'on y menos servicio de deuda, en magnitudes del orden
de varios puntos del PIB.

\subsection{Por qu\'e la diferencia no implica ``mejor''}

GPT-4o-mini termina con menor pobreza y menor deuda, pero el
simulador no representa fielmente el costo de subfinanciar el
servicio de deuda en una econom\'ia con alta dependencia de
financiamiento externo. Un modelo m\'as conservador como Claude
podr\'ia estar viendo un riesgo que el sim no penaliza. La
conclusi\'on correcta no es ``GPT $>$ Claude para Guatemala'', sino
``el sim revela las prioridades que cada modelo trae, y esas
prioridades difieren''.

\subsection{Limitaciones}

(1)~$N{=}1$ corrida por modelo en cada \emph{seed}; los resultados
son estables entre las dos corridas pero no constituyen un test
estad\'istico riguroso. El siguiente paso es $\geq 30$ \emph{seeds}
$\times$ 16 turnos, aproximadamente USD~15 de costo de API y
$\sim 1.5$ horas. (2)~Horizonte corto (8 turnos $\approx$ 2 a\~nos);
no medimos efectos compuestos de largo plazo. (3)~Din\'amica del
mundo simplificada, sin ciclo electoral ni \emph{shocks}
geopol\'iticos end\'ogenos. (4)~Memoria presidencial limitada: cada
llamada al LLM es independiente, sin \emph{carry-over} de mensajes;
la \'unica memoria \emph{cross-turno} est\'a en
\texttt{state.memoria\_presidencial} serializado. (5)~Solo dos
modelos frontera; falta Gemini, Llama 3.1 405B, DeepSeek.
(6)~Calibraci\'on del estado inicial contra datos p\'ublicos pero no
auditada por economistas guatemaltecos: las magnitudes relativas son
m\'as confiables que los niveles.

\subsection{Trabajo Futuro}

Barrido de $\geq 30$ \emph{seeds} con \emph{intra-class correlation}
para distinguir se\~nal de modelo de ruido del \emph{sampler}. Adici\'on
de Gemini, Llama 3.1, DeepSeek-V3 para test cross-vendor de la
hip\'otesis $H_{TC}$ (\emph{transfer cultural}). Replicaci\'on en
Honduras, Chile y Bolivia: \emph{datasets} de calibraci\'on locales,
metodolog\'ia id\'entica---primer paso de la agenda ``AI Safety from
the Global South''. Estudio de ablaci\'on del \emph{system prompt}:
?`cu\'anto del comportamiento ``constitucional'' es atribuible al
\emph{prompt} vs.\ al modelo?

% =====================================================================
\section{Conclusi\'on}

Construimos un \emph{testbed} bayesiano de auditor\'ia de
LLM-como-decisor calibrado contra Guatemala 2026, midiendo tanto qu\'e
le pasa al pa\'is como qu\'e revela el LLM sobre sus propias
prioridades. Bajo id\'enticos \emph{shocks} e id\'entico contexto,
\textbf{Claude Haiku 4.5 y GPT-4o-mini gobiernan distinto, y la
diferencia es medible y reproducible}: Claude prioriza el servicio
de deuda (19\,\% del presupuesto) y la protecci\'on social;
GPT-4o-mini prioriza salud y educaci\'on ($\approx$\,18\,\% cada una)
con m\'inima asignaci\'on a deuda (5\,\%). En \emph{outcomes},
GPT-4o-mini termina con menor pobreza, mayor aprobaci\'on y menor
deuda. Ambos se desv\'ian del \emph{baseline} humano MINFIN 2024 en
direcciones opuestas.

La pregunta que el paper plantea---y deja expl\'icitamente
abierta---es: si la elecci\'on entre dos modelos frontera implica
varios puntos porcentuales del PIB en asignaci\'on presupuestaria,
\emph{?`qui\'en decide qu\'e modelo se usa cuando un LLM se incorpora
al ciclo de pol\'itica p\'ublica?} El m\'etodo aqu\'i propuesto da una
herramienta concreta para responderla \emph{ex-ante}, calibrada al
pa\'is donde se va a desplegar.

% =====================================================================
\section*{Reproducibilidad}

C\'odigo, datos y figuras: repositorio \texttt{guatemala-sim} (commit
\texttt{121fe35}+). Run principal:
\texttt{20260419\_224225\_836edc}, seed~11. Run de validaci\'on:
\texttt{20260419\_222128\_9d4e55}. Suite de tests: 254 \emph{tests}
\emph{offline} (\texttt{python -m pytest}).

% =====================================================================
\begin{thebibliography}{00}

\bibitem{ramachandran2007}
D.~Ramachandran and E.~Amir, ``Bayesian inverse reinforcement learning,''
in \emph{Proc.\ IJCAI}, 2007, pp.\ 2586--2591.

\bibitem{hadfield2017}
D.~Hadfield-Menell, S.~Milli, P.~Abbeel, S.~Russell, and A.~Dragan,
``Inverse reward design,'' in \emph{Adv.\ Neural Inf.\ Process.\ Syst.\
(NeurIPS)}, 2017.

\bibitem{ziebart2008}
B.~D.~Ziebart, A.~Maas, J.~A.~Bagnell, and A.~K.~Dey, ``Maximum
entropy inverse reinforcement learning,'' in \emph{Proc.\ AAAI}, 2008.

\bibitem{lanham2023}
T.~Lanham et al., ``Measuring faithfulness in chain-of-thought
reasoning,'' \emph{arXiv:2307.13702}, 2023.

\bibitem{park2023}
J.~S.~Park et al., ``Generative agents: Interactive simulacra of human
behavior,'' in \emph{Proc.\ UIST}, 2023.

\bibitem{liu2023}
X.~Liu et al., ``AgentBench: Evaluating LLMs as agents,''
\emph{arXiv:2308.03688}, 2023.

\bibitem{pan2023}
A.~Pan et al., ``Do the rewards justify the means? Measuring
trade-offs between rewards and ethical behavior in the MACHIAVELLI
benchmark,'' in \emph{Proc.\ ICML}, 2023.

\bibitem{perez2022}
E.~Perez et al., ``Discovering language model behaviors with
model-written evaluations,'' \emph{arXiv:2212.09251}, 2022.

\bibitem{willard2023}
B.~T.~Willard and R.~Louf, ``Outlines: Efficient structured
generation,'' \emph{arXiv:2307.09702}, 2023.

\bibitem{anthropic_rsp}
Anthropic, ``Responsible Scaling Policy,'' \emph{Anthropic technical
report}, 2024.

\bibitem{deepmind_fsf}
Google DeepMind, ``Frontier Safety Framework,'' \emph{DeepMind
technical report}, 2024.

\bibitem{nist_aimrf}
NIST, ``AI Risk Management Framework (AI RMF~1.0),'' NIST AI 100-1,
2023.

\bibitem{aisi_inspect}
UK AISI, ``Inspect: An open-source framework for large language model
evaluations,'' 2024.

\bibitem{wb_wdi}
World Bank, ``World Development Indicators---Guatemala,'' 2024.

\bibitem{banguat}
Banco de Guatemala, ``Estad\'isticas macroecon\'omicas,'' 2024--2026.

\bibitem{minfin}
Ministerio de Finanzas P\'ublicas de Guatemala, ``Liquidaci\'on
presupuestaria,'' 2024.

\end{thebibliography}

\end{document}
